\section{Unblocked Algorithms}
\label{sec:lu-unblocked}

\subsection{Six Loop Orders}

LU 분해의 산술은 세 겹의 루프 하나로 적힌다. 피벗을 잠시 잊으면
\[
\text{for } k,\ \text{for } i>k,\ \text{for } j>k:\quad a_{ij}\gets a_{ij}-\ell_{ik}a_{kj},\qquad \ell_{ik}=a_{ik}/a_{kk}
\]
이고 $(k,i,j)$의 여섯 가지 순서가 모두 같은 산술을 한다. 원소 $a_{ij}$가 받는 빼기는 언제나 $k=1,2,\dots$의 순서로 일어나므로 반올림까지 같은 결과가 나온다. 다른 것은 메모리를 읽고 쓰는 모양이다. \cref{fig:lu-looking}는 세 가지 대표 꼴을 보인다.
\begin{itemize}
\item \term{오른쪽 보기}{right-looking}($kij$, $kji$): $k$번째 단계에서 승수 열을 만들자마자 오른쪽 아래 남은 행렬 전체를 갱신한다. 갱신이 랭크 1 행렬 $\ell_ku_k^{\mathsf T}$를 빼는 것이어서 바깥곱 꼴이라고도 한다.
\item \term{왼쪽 보기}{left-looking}($jki$, $jik$): $j$열을 만들 차례가 되어서야 왼쪽에 이미 만든 열들로 $j$열 하나를 갱신한다. 한 단계에 쓰는 것은 한 열뿐이다.
\item \term{행 지향 꼴}{row-oriented}($ikj$, $ijk$): $i$행을 만들 차례에 위에 이미 만든 행들로 $L$과 $U$의 $i$행을 한꺼번에 만든다. 왼쪽 보기를 전치한 꼴이다.
\item \term{Crout 꼴}{Crout form}: $k$번째 단계에 $U$의 $k$행과 $L$의 $k$열을 이미 만든 부분과의 내적으로 만든다. 행 지향 $ijk$의 내적과 열 지향 $jik$의 내적을 번갈아 쓴 것이다.
\end{itemize}

\begin{figure}[t]
\centering
\begin{tikzpicture}[x=0.36cm,y=0.36cm,font=\small]
\foreach \xo/\name in {0/오른쪽 보기,9/왼쪽 보기,18/Crout 꼴}{
  \draw (\xo,0) rectangle (\xo+7,7);
  \node at (\xo+3.5,-1) {\name};
}
% right-looking: done L cols + U rows in light gray; current col/row dark; trailing updated hatched
\fill[gray!15] (0,0) rectangle (3,7); \fill[gray!15] (3,4) rectangle (7,7);
\fill[gray!55] (3,0) rectangle (4,4); \fill[gray!55] (3,3) rectangle (7,4);
\fill[pattern=north east lines] (4,0) rectangle (7,3); \draw (4,0) rectangle (7,3);
% left-looking: done cols light; read L part of done cols; column j written
\fill[gray!15] (9,0) rectangle (12,7);
\fill[gray!55] (12,0) rectangle (13,7);
\draw[pattern=north east lines] (12,0) rectangle (13,7);
% crout: row k of U and column k of L written; done rows/cols read
\fill[gray!15] (18,0) rectangle (21,7); \fill[gray!15] (21,4) rectangle (25,7);
\fill[gray!55] (21,0) rectangle (22,4); \fill[gray!55] (21,3) rectangle (25,4);
\end{tikzpicture}
\caption{세 가지 루프 꼴이 한 단계에 읽고 쓰는 곳. 연한 회색은 이미 끝난 $L$과 $U$, 진한 회색은 이 단계에 완성되는 부분, 빗금은 이 단계에 갱신되는 부분이다. 오른쪽 보기는 남은 행렬 전체를 매 단계 고치고, 왼쪽 보기는 한 열만 고치며, Crout 꼴은 남은 행렬을 건드리지 않는다.}
\label{fig:lu-looking}
\end{figure}

열 우선(column-major) 배열에서는 같은 열의 원소가 메모리에 붙어 있으므로 가장 안쪽 루프가 $i$(행)를 돌아야 연속한 메모리를 훑는다. 그래서 열 우선에 맞는 꼴은 $kji$(오른쪽 보기)와 $jki$(왼쪽 보기)이고 LAPACK의 \code{getf2}는 $kji$를 쓴다. 부분 피벗을 넣으면 다음과 같다.

\begin{algorithm}[t]
\begin{algorithmic}[1]
\Require $m\times n$ 행렬 $A$ (열 우선, 그 자리에 덮어씀)
\Ensure $A$의 대각 아래에 $L$의 승수, 대각과 위에 $U$; 교환 기록 $\mathit{ipiv}$; 정확히 $0$인 첫 피벗의 번호 $\mathit{info}$
\State $\mathit{info}\gets0$
\For{$k=1,\dots,\min(m,n)$}
  \State $p\gets\arg\max_{k\le i\le m}|a_{ik}|$ \Comment{같은 값이면 가장 앞의 $i$}
  \State $\mathit{ipiv}_k\gets p$
  \If{$a_{pk}=0$}
    \State \textbf{if} $\mathit{info}=0$ \textbf{then} $\mathit{info}\gets k$ \Comment{특이: 이 열은 소거하지 않는다}
    \State \textbf{continue}
  \EndIf
  \State 행 $k$와 행 $p$를 열 $1,\dots,n$ 전체에서 바꾼다
  \State $a_{k+1:m,\,k}\gets a_{k+1:m,\,k}/a_{kk}$ \Comment{승수}
  \For{$j=k+1,\dots,n$} \Comment{랭크 1 갱신, 열 하나씩}
    \For{$i=k+1,\dots,m$}
      \State $a_{ij}\gets a_{ij}-a_{ik}\,a_{kj}$
    \EndFor
  \EndFor
\EndFor
\end{algorithmic}
\caption{부분 피벗을 쓰는 오른쪽 보기 LU ($kji$ 꼴, LAPACK \code{getf2}의 산술)}
\label{alg:lu-right}
\end{algorithm}

\cref{alg:lu-right}의 줄마다 할 일을 짚는다.
\begin{itemize}
\item 피벗 찾기(3행)는 $k$열의 대각부터 아래를 훑는다. 절댓값이 같은 후보가 여럿이면 처음 것을 고르는 규칙은 결과를 정하는 규칙이기도 하다. \cref{prop:lu-growth-bound}의 $W_n$이 교환 없이 소거되는 것이 이 규칙 때문이다.
\item 피벗이 정확히 $0$이면(5행) 그 열의 대각 아래도 모두 $0$이므로 소거할 것이 없다. 알고리즘은 멈추지 않고 첫 $0$ 피벗의 번호만 기록한다. 그래서 특이 행렬도 \cref{sec:lu-factorization}의 뜻에서 분해가 끝까지 나오고, 그것을 에러로 볼지는 부르는 쪽이 정한다.
\item 행 교환(9행)은 이미 끝난 왼쪽 열들(승수)까지 바꾼다. \cref{thm:lu-partial-pivoting}의 증명에서 승수를 뒤의 교환으로 섞어야 했던 것이 바로 이 일이다. 이 덕분에 끝난 뒤 $A$의 대각 아래가 그대로 $PA=LU$의 $L$이 된다.
\item 승수(10행)는 피벗으로 나눈다. LAPACK의 \code{getf2}는 $1/a_{kk}$를 한 번 구해 곱하는데, $|a_{kk}|$가 가장 작은 정규 수보다 작으면 역수가 넘칠 수 있으므로 그때만 원소마다 나눈다. 곱셈이 나눗셈보다 훨씬 빠르기 때문에 고른 길이며, 역수를 곱한 결과는 나눈 결과와 반올림 한 번만큼 다를 수 있다.
\item 갱신(11--15행)은 열마다 \code{axpy} 한 번이다. LAPACK은 이 이중 루프를 BLAS의 랭크 1 갱신 \code{ger} 하나로 부른다.
\end{itemize}

이 알고리즘의 \term{루프 불변식}{loop invariant}을 적어 두면 뒤의 블록 알고리즘과 재귀 알고리즘이 같은 결과를 낸다는 것을 보이기 쉽다. $k$번째 반복이 끝난 뒤 지금까지의 교환을 모은 순열을 $P^{(k)}=P_k\cdots P_1$이라 하면 배열은
\begin{equation}
\label{eq:lu-invariant}
\begin{pmatrix}L_{11}\backslash U_{11}&U_{12}\\L_{21}&S\end{pmatrix},\qquad
P^{(k)}A=\begin{pmatrix}L_{11}&0\\L_{21}&I\end{pmatrix}\begin{pmatrix}U_{11}&U_{12}\\0&S\end{pmatrix}
\end{equation}
을 담고 있다. $L_{11}\backslash U_{11}$은 $k\times k$ 블록의 대각 아래에 $L_{11}$, 대각과 위에 $U_{11}$을 겹쳐 적은 것이고 $S$는 $P^{(k)}A$의 왼쪽 위 $k\times k$ 블록에 대한 Schur 보수이다. $k=0$이면 $S=A$이다. 반복 하나가 불변식을 지키는 것은 \eqref{eq:lu-schur-step}을 $S$에 쓰고 \cref{thm:lu-partial-pivoting}의 증명처럼 새 교환을 왼쪽 열들의 승수에도 적용한 결과이다. $k=\min(m,n)$이면 $S$가 비거나 $0$이어서 \eqref{eq:lu-invariant}가 곧 $PA=LU$이다.

\begin{algorithm}[t]
\begin{algorithmic}[1]
\Require $m\times n$ 행렬 $A$ (열 우선, 그 자리에 덮어씀)
\Ensure \cref{alg:lu-right}과 같은 $L$, $U$, $\mathit{ipiv}$
\For{$j=1,\dots,n$}
  \State 열 $j$에 교환 $\mathit{ipiv}_1,\dots,\mathit{ipiv}_{\min(j-1,m)}$을 차례로 적용 \Comment{미뤄 둔 교환}
  \State $k\gets\min(j-1,m)$
  \State $a_{1:k,\,j}\gets L_{1:k,1:k}^{-1}a_{1:k,\,j}$ \Comment{\code{trsv}: 단위 하삼각, $U$의 열 $j$ 윗부분}
  \State $a_{k+1:m,\,j}\gets a_{k+1:m,\,j}-A_{k+1:m,\,1:k}\,a_{1:k,\,j}$ \Comment{\code{gemv}: 왼쪽 열들의 기여}
  \If{$j\le m$}
    \State $p\gets\arg\max_{j\le i\le m}|a_{ij}|$; $\mathit{ipiv}_j\gets p$
    \State 행 $j$와 행 $p$를 열 $1,\dots,j$에서 바꾼다 \Comment{오른쪽 열은 그 차례에}
    \State \textbf{if} $a_{jj}\ne0$ \textbf{then} $a_{j+1:m,\,j}\gets a_{j+1:m,\,j}/a_{jj}$
  \EndIf
\EndFor
\end{algorithmic}
\caption{부분 피벗을 쓰는 왼쪽 보기 LU ($jki$ 꼴)}
\label{alg:lu-left}
\end{algorithm}

\cref{alg:lu-left}은 열 $j$를 만들 차례가 될 때까지 그 열을 건드리지 않는다. 그래서 그 사이에 일어난 교환을 열 $j$에 미뤄 두었다가 2행에서 한꺼번에 적용한다. 4행의 삼각 풀이는 \cref{prop:lu-blocked-same}의 증명에서 본 $u_{kj}=a_{kj}-\sum_{i<k}\ell_{ki}u_{ij}$를 열 하나에 대해 하는 것이고, 5행은 그 열의 대각 아래에 왼쪽 모든 열의 갱신을 한꺼번에 몰아 준다. 산술은 \cref{alg:lu-right}과 같고 메모리를 쓰는 곳은 매 단계 열 하나뿐이다.

\subsection{Operation Count}

$k$번째 단계는 나눗셈 $m-k$번과 갱신 $(m-k)(n-k)$개 원소에 곱셈과 뺄셈 하나씩이다. 정사각이면
\[
\sum_{k=1}^{n-1}\bigl[(n-k)+2(n-k)^2\bigr]=\frac{n(n-1)}2+\frac{(n-1)n(2n-1)}3=\frac23n^3-\frac12n^2-\frac16n
\]
플롭이다. 이 $(2/3)n^3$이 LU의 비용으로 늘 인용되는 수이다.

\subsection{Memory Traffic}

메모리의 비용을 세려면 모형이 필요하다. 산술 장치 곁에 $M$개의 수를 담는 빠른 메모리(캐시)가 있고 그 바깥에 큰 느린 메모리가 있다고 하자. 알고리즘의 \term{통신량}{communication volume}은 두 메모리 사이를 오간 수의 개수이다. 행렬이 캐시보다 훨씬 크다고 할 때 세 루프 꼴의 통신량은 다음과 같다.
\begin{itemize}
\item 오른쪽 보기: $k$번째 단계가 남은 $(m-k)(n-k)$개를 한 번씩 읽고 쓰므로 정사각에서 $\sum_k2(n-k)^2\approx(2/3)n^3$개이다. 플롭 수도 $(2/3)n^3$이므로 수 하나를 옮길 때마다 플롭 하나를 한다.
\item 왼쪽 보기: $j$번째 단계가 왼쪽 $j-1$열의 $L$(삼각 부분 $j^2/2$개와 아래 부분 $(n-j)j$개)을 읽고 열 하나를 쓴다. 읽기는 $\sum_j(jn-j^2/2)\approx n^3/3$개, 쓰기는 $n^2$개이다. 통신량이 오른쪽 보기의 절반이고 대부분 읽기여서 쓰기가 비싼 저장 장치(디스크 밖의 분해)에서는 이 꼴이 유리하다.
\item 블록 오른쪽 보기(폭 $b$): 단계마다 남은 행렬 $C$를 한 번 읽고 쓰고 패널과 $U_{12}$를 읽는다. 단계가 $n/b$번이므로 통신량은 $\sum_j2(n-j)^2\approx(2/3)n^3/b$개이고 수 하나에 플롭 $b$개이다.
\end{itemize}
블록의 이득은 $b$가 클수록 크지만 \code{gemm}이 정말 수 하나를 $b$번 쓰려면 $C$의 $r\times r$ 조각을 고치는 동안 $A_{21}$의 $r\times b$ 조각과 $U_{12}$의 $b\times r$ 조각, 그리고 $C$의 조각이 캐시에 함께 있어야 한다. $r=b$로 잡으면 $3b^2\le M$이어야 하므로 $b\approx\sqrt{M/3}$이 한계이고 통신량은 $O(n^3/\sqrt M)$이다.

이것이 최선이라는 하한이 있다. Hong과 Kung은 고전적인 행렬 곱이 $\Omega(n^3/\sqrt M)$개의 수를 옮겨야 한다는 것을 보였고\cite{hong1981io}, Irony, Toledo, Tiskin의 증명\cite{irony2004communication}을 Ballard, Demmel, Holtz, Schwartz가 LU를 포함한 직접법 전체로 넓혔다\cite{ballard2011minimizing}. 증명의 줄거리는 이렇다. 계산을 통신 $M$번씩의 구간으로 자른다. 한 구간 안에서 캐시에 있을 수 있는 수는 많아야 $2M$개(들어 있던 $M$개와 새로 들어온 $M$개)이다. 갱신 $a_{ij}\gets a_{ij}-\ell_{ik}u_{kj}$ 하나는 세 첨자 쌍 $(i,j)$, $(i,k)$, $(k,j)$를 쓰므로, 한 구간에서 할 수 있는 갱신의 집합 $V$는 세 좌표 평면으로의 사영이 각각 $2M$개 이하인 격자점 집합이다. Loomis--Whitney 부등식 $|V|\le\sqrt{|V_{ij}||V_{ik}||V_{kj}|}$에 따라 $|V|\le(2M)^{3/2}$이다. 갱신은 모두 약 $n^3/3$개이므로 구간이 $n^3/(3(2M)^{3/2})$개 넘게 필요하고 통신량은 $\Omega(n^3/\sqrt M)$이다. 블록 LU는 이 하한을 상수배까지 이룬다.

이 모형으로 \cref{fig:lu-gflops}의 맨 아래 곡선을 설명할 수 있다. 오른쪽 보기의 갱신은 수 하나(8바이트)를 읽고 쓰는 동안 플롭 둘을 하므로 바이트당 $1/8$ 플롭이다. 이 기계에서 한 코어로 잰 메모리 대역폭(STREAM의 triad 꼴, \code{figures/48/extra.c})은 초당 $131$GB이고 여기에 $1/8$을 곱한 초당 $16\times10^9$ 플롭이 블록 없는 LU의 천장이다. 측정한 $12\times10^9$ 플롭은 그 아래에 있다. 이렇게 처리 속도를 산술 강도와 대역폭의 곱으로 묶는 그림을 \term{지붕선 모형}{roofline model}이라 한다.

재귀 LU의 통신량은 점화식으로 센다. $m\times n$ 패널($m\ge n$)의 통신량을 $Q(m,n)$이라 하면 패널이 캐시에 들어갈 때($mn\le M$) $Q=O(mn)$이고 그 밖에는
\[
Q(m,n)\le Q(m,n/2)+Q(m-n/2,\,n/2)+O\Bigl(\frac{mn^2}{\sqrt M}+mn\Bigr)
\]
이다. 마지막 항은 크기 $m\times n/2\times n/2$의 \code{trsm}과 \code{gemm}을 캐시를 모르는 방법으로 했을 때의 통신량이다. 재귀의 $\ell$번째 층에는 폭 $n/2^\ell$의 문제가 $2^\ell$개 있어 첫 항의 합이 $mn^2/(2^\ell\sqrt M)$이고 층을 따라 등비로 줄어든다. 둘째 항은 층마다 $O(mn)$이고 층은 $\log_2n$개이다. 그래서 $Q(m,n)=O(mn^2/\sqrt M+mn\log n)$이고 $n$이 $\sqrt M\log n$보다 크면 첫 항이 지배해 하한과 같은 차수가 된다\cite{toledo1997recursive}. 블록 폭이 어디에도 나타나지 않는다는 것이 캐시를 모르는 알고리즘의 뜻이다.

\section{Blocked and Recursive LU}
\label{sec:lu-blocked}

\subsection{The BLAS Levels}

BLAS(Basic Linear Algebra Subprograms)는 선형대수의 기본 연산을 셋으로 나눈다\cite{lawson1979blas1,dongarra1990blas3}. 1단계는 벡터 연산(\code{axpy}, \code{dot}; $O(n)$ 데이터에 $O(n)$ 연산), 2단계는 행렬–벡터 연산(\code{gemv}, \code{ger}, \code{trsv}; $O(n^2)$ 데이터에 $O(n^2)$ 연산), 3단계는 행렬–행렬 연산(\code{gemm}, \code{trsm}; $O(n^2)$ 데이터에 $O(n^3)$ 연산)이다. 데이터 하나를 메모리에서 가져와 몇 번의 연산에 쓰는지(산술 강도)가 1, 2단계에서는 상수이고 3단계에서만 $n$에 비례한다. 3단계 연산은 행렬을 캐시와 레지스터에 맞는 블록으로 잘라 블록 하나를 여러 번 재사용할 수 있으므로 산술 장치의 최고 속도에 가깝게 돈다. 빠른 LU란 산술의 대부분을 \code{gemm}으로 바꾼 LU이다.

\subsection{Block Right-Looking LU}

$A$의 앞 $b$열을 \term{패널}{panel}로 떼어 블록으로 나눈다.
\[
A=\begin{pmatrix}A_{11}&A_{12}\\A_{21}&A_{22}\end{pmatrix},\qquad A_{11}\in\mathbb C^{b\times b}.
\]
패널 $\binom{A_{11}}{A_{21}}$을 부분 피벗으로 분해해 $P_1\binom{A_{11}}{A_{21}}=\binom{L_{11}}{L_{21}}U_{11}$을 얻었다고 하자. 같은 교환을 오른쪽 열에도 적용하고 교환된 블록을 다시 $A_{ij}$로 쓰면 $P_1A$는 다음처럼 나뉜다.
\[
\begin{pmatrix}A_{11}&A_{12}\\A_{21}&A_{22}\end{pmatrix}
=\begin{pmatrix}L_{11}&0\\L_{21}&I\end{pmatrix}
\begin{pmatrix}U_{11}&U_{12}\\0&S\end{pmatrix},\qquad
U_{12}=L_{11}^{-1}A_{12},\quad S=A_{22}-L_{21}U_{12}
\]
이다. 식 \eqref{eq:lu-schur-step}을 블록으로 바꾼 것이며 남은 일은 $S$를 같은 방법으로 분해하는 것이다(\cref{fig:lu-blocked}). $U_{12}$를 구하는 것은 단위 하삼각 행렬로 여러 우변을 푸는 \code{trsm}이고, $S$의 갱신은 $(n-b)\times b$와 $b\times(n-b)$의 곱인 \code{gemm}이다.

\begin{figure}[t]
\centering
\begin{tikzpicture}[x=0.42cm,y=0.42cm,font=\small]
\draw (0,0) rectangle (9,9);
\fill[gray!15] (0,0) rectangle (2,9); \fill[gray!15] (2,7) rectangle (9,9);
\fill[gray!55] (2,0) rectangle (4,7);
\fill[gray!35] (4,5) rectangle (9,7);
\fill[pattern=north east lines] (4,0) rectangle (9,5); \draw (4,0) rectangle (9,5);
\draw (2,0) -- (2,9) (4,0) -- (4,7) (2,7) -- (9,7) (2,5) -- (9,5);
\node at (1,4.5) {끝남};
\node[fill=white,inner sep=1pt] at (3,6) {$A_{11}$};
\node[fill=white,inner sep=1pt] at (3,2.5) {$A_{21}$};
\node at (6.5,6) {$A_{12}$};
\node[fill=white,inner sep=1pt] at (6.5,2.5) {$A_{22}$};
\begin{scope}[xshift=4.6cm]
\node[anchor=west,align=left,font=\small] at (0,6.8) {(i) 패널 분해 \code{getrf2}: 진한 회색};
\node[anchor=west,align=left,font=\small] at (0,5.3) {(ii) 교환을 좌우 열에 \code{laswp}};
\node[anchor=west,align=left,font=\small] at (0,3.8) {(iii) $U_{12}=L_{11}^{-1}A_{12}$ \code{trsm}: 중간 회색};
\node[anchor=west,align=left,font=\small] at (0,2.3) {(iv) $A_{22}\gets A_{22}-L_{21}U_{12}$ \code{gemm}: 빗금};
\end{scope}
\end{tikzpicture}
\caption{블록 오른쪽 보기 LU의 한 단계. 패널은 폭 $b$의 열 묶음이고 이 단계의 산술 대부분은 빗금 친 부분의 \code{gemm}이다.}
\label{fig:lu-blocked}
\end{figure}

\begin{algorithm}[t]
\begin{algorithmic}[1]
\Require $m\times n$ 행렬 $A$, 블록 폭 $b$ (LAPACK 참조 구현의 기본값 $64$)
\Ensure \cref{alg:lu-right}과 같은 꼴의 $L$, $U$, $\mathit{ipiv}$, $\mathit{info}$
\If{$b\le1$ \textbf{or} $b\ge\min(m,n)$}
  \State \Return \Call{getrf2}{$A$} \Comment{작으면 블록 없이}
\EndIf
\For{$j=1,\,1+b,\,1+2b,\dots\le\min(m,n)$}
  \State $j_b\gets\min(\min(m,n)-j+1,\ b)$
  \State \Call{getrf2}{$A_{j:m,\,j:j+j_b-1}$} \Comment{패널; 피벗 번호는 패널 안에서 센다}
  \State $\mathit{ipiv}_{j:j+j_b-1}\gets\mathit{ipiv}_{j:j+j_b-1}+j-1$; $\mathit{info}$도 같이 옮긴다
  \State 행 교환 $\mathit{ipiv}_{j:j+j_b-1}$을 열 $1,\dots,j-1$에 적용 \Comment{\code{laswp}: 끝난 $L$}
  \If{$j+j_b\le n$}
    \State 같은 교환을 열 $j+j_b,\dots,n$에 적용 \Comment{\code{laswp}}
    \State $A_{j:j+j_b-1,\,j+j_b:n}\gets L_{11}^{-1}A_{j:j+j_b-1,\,j+j_b:n}$ \Comment{\code{trsm}, $L_{11}$은 단위 하삼각}
    \If{$j+j_b\le m$}
      \State $A_{j+j_b:m,\,j+j_b:n}\gets A_{j+j_b:m,\,j+j_b:n}-A_{j+j_b:m,\,j:j+j_b-1}A_{j:j+j_b-1,\,j+j_b:n}$ \Comment{\code{gemm}}
    \EndIf
  \EndIf
\EndFor
\end{algorithmic}
\caption{블록 오른쪽 보기 LU (LAPACK \code{dgetrf}의 구조를 그대로 옮김)}
\label{alg:lu-blocked}
\end{algorithm}

\cref{alg:lu-blocked}은 LAPACK 참조 구현 \code{dgetrf}의 순서를 그대로 옮긴 것이다(\cite{anderson1999lapack}, 참조 구현의 소스로 확인). 몇 가지를 짚는다.
\begin{itemize}
\item 패널 분해는 $(m-j+1)\times j_b$의 길쭉한 행렬에 대한 LU이고 피벗을 패널 높이 전체에서 찾는다. 그래서 블록 알고리즘도 피벗 규칙은 \cref{alg:lu-right}과 같은 부분 피벗이며, 반올림을 빼면 같은 $P$, $L$, $U$를 낸다.
\item 패널이 돌려준 피벗 번호는 패널의 첫 행을 $1$로 센 것이므로 7행에서 $j-1$을 더해 전체 행 번호로 바꾼다.
\item 교환은 두 번에 나누어 적용한다. 패널 안의 열은 패널 분해가 이미 바꿨고, 왼쪽의 끝난 열(8행)과 오른쪽의 아직 손대지 않은 열(10행)을 따로 바꾼다.
\item 산술량을 세 부분으로 나누면 패널 분해가 단계마다 약 $(n-j)j_b^2$, \code{trsm}이 약 $j_b^2(n-j)$이고 나머지가 모두 \code{gemm}이다. 단계들을 더하면 \code{gemm}이 아닌 부분이 약 $n^2b$이므로 \code{gemm}의 몫은
\[
1-\frac{n^2b}{2n^3/3}=1-\frac{3b}{2n}
\]
이다. $n=1024$, $b=64$이면 $91\%$이다. $b$를 키우면 이 몫이 줄고, 줄이면 \code{gemm} 하나하나가 납작해져(내적 길이 $b$) 재사용이 줄어든다. $b$는 이 둘 사이의 타협이고 하드웨어마다 다르다. 참조 구현은 \code{ilaenv}가 돌려주는 $64$를 쓴다.
\end{itemize}

\begin{proposition}
\label{prop:lu-blocked-same}
정확한 산술에서 \cref{alg:lu-blocked}(패널을 \cref{alg:lu-right}로 분해할 때)은 같은 동점 규칙의 \cref{alg:lu-right}과 같은 $\mathit{ipiv}$, $L$, $U$를 낸다.
\end{proposition}

\begin{proof}
첫 패널에 대해 보이면 나머지는 Schur 보수 $S$에 대한 귀납이다. \cref{alg:lu-right}의 $k$번째 단계($k\le b$)가 열 $j$를 바꾸는 방법은 행 교환과 $a_{ij}\gets a_{ij}-a_{ik}a_{kj}$뿐이고, 여기에 쓰이는 값은 열 $k$와 열 $j$의 원소뿐이다. 그러므로 처음 $b$번의 단계 동안 패널의 열 $1,\dots,b$가 겪는 일은 패널 밖의 열과 무관하고, 패널만 따로 분해한 \cref{alg:lu-blocked}의 6행과 정확히 같다. 피벗은 패널의 열에서만 고르므로 처음 $b$개의 피벗과 승수도 같다. 패널 밖의 열 $j>b$는 같은 교환을 받고 행 $k\le b$에서 $u_{kj}=a_{kj}-\sum_{i<k}\ell_{ki}u_{ij}$가 되는데, 이것은 단위 하삼각 $L_{11}$로 전진 대입한 $U_{12}=L_{11}^{-1}A_{12}$의 열 $j$이다. 행 $i>b$에서는 $a_{ij}-\sum_{k\le b}\ell_{ik}u_{kj}$, 곧 $A_{22}-L_{21}U_{12}$의 원소가 된다. 교환이 모두 행 $k$ 이상에서 일어나므로 교환을 먼저 다 하고 갱신해도 같다(\cref{thm:lu-partial-pivoting}의 증명). 이것이 \cref{alg:lu-blocked}의 8--13행이다.
\end{proof}

반올림이 있으면 두 알고리즘은 빼는 순서가 달라(\code{gemm}은 $k$에 대한 합을 블록으로 묶는다) 원소의 끝 비트가 다를 수 있고, 피벗 후보의 크기가 거의 같으면 다른 행을 고를 수도 있다. \cref{thm:lu-backward}의 상한은 빼는 순서와 상관없으므로 두 결과 모두 같은 오차 해석을 따른다.

\cref{fig:lu-gflops}는 이 차이를 보인다. 직접 짠 \cref{alg:lu-blocked}(패널은 \cref{alg:lu-right}, \code{trsm}과 \code{gemm}은 Accelerate)이 블록 없는 루프보다 $n=1536$에서 일곱 배 빠르다. 같은 그림에서 $n=256$, $512$, $1024$의 블록 곡선이 움푹 꺼지는 것은 우연이 아니다. 이때 열 하나의 길이가 $2^{k}$의 배수여서 같은 행의 원소들이 캐시의 같은 묶음(set)에 몰린다. leading dimension만 $512$에서 $520$으로 늘려 같은 분해를 하면 초당 $36\times10^9$에서 $85\times10^9$ 플롭으로, $1024$에서 $1032$로 늘리면 $59\times10^9$에서 $122\times10^9$ 플롭으로 뛴다. 같은 산술, 같은 메모리 양인데 주소의 모양만 바뀐 결과이다.

\begin{figure}[t]
\centering
\begin{tikzpicture}
\begin{semilogxaxis}[width=0.86\textwidth,height=0.55\textwidth,
  xlabel={$n$},ylabel={$10^9$ 플롭/초},grid=major,grid style={gray!25},
  legend pos=north west,legend style={font=\small,draw=none,fill=none},
  xmin=28,xmax=2300,ymin=0,ymax=300,log basis x=2]
\addplot[black,mark=*,mark size=1.2pt] table[x=n,y=getrf]{figures/48/flops.dat};
\addlegendentry{Accelerate \code{dgetrf}}
\addplot[black,dashed,mark=square*,mark size=1.2pt] table[x=n,y=blocked]{figures/48/flops.dat};
\addlegendentry{\cref{alg:lu-blocked} ($b=64$)}
\addplot[black!60,mark=triangle*,mark size=1.4pt] table[x=n,y=getf2]{figures/48/flops.dat};
\addlegendentry{Accelerate \code{dgetf2}}
\addplot[black!60,dotted,thick,mark=o,mark size=1.2pt] table[x=n,y=unblocked]{figures/48/flops.dat};
\addlegendentry{\cref{alg:lu-right} (C 루프)}
\end{semilogxaxis}
\end{tikzpicture}
\caption{정사각 LU의 속도(Apple M4 Pro, 플롭 수는 $(2/3)n^3$으로 셈). 측정 코드는 \code{figures/48/bench.c}이다. \code{dgetrf}는 $n\ge768$에서 Accelerate가 안에서 여러 코어를 쓴다.}
\label{fig:lu-gflops}
\end{figure}

\subsection{Recursive LU}

패널을 \cref{alg:lu-right}처럼 블록 없이 분해하면(LAPACK 3.5까지의 \code{dgetrf}는 패널을 \code{dgetf2}로 분해했다) 패널 분해는 여전히 2단계 BLAS이다. 패널 분해의 산술은 전체의 약 $3b/(4n)$이지만 메모리에 묶인 속도로 돌기 때문에 시간으로는 그보다 큰 몫을 차지한다. Toledo\cite{toledo1997recursive}와 Gustavson\cite{gustavson1997recursion}은 패널을 다시 반으로 나누는 재귀로 이것을 없앴다. 열을 반으로 나누어 왼쪽 반을 재귀로 분해하고, 오른쪽 반을 \code{trsm}과 \code{gemm}으로 갱신한 뒤, 남은 아래 오른쪽을 재귀로 분해한다. LAPACK 3.6부터 패널 분해 \code{getrf2}가 이 재귀이다.

\begin{algorithm}[t]
\begin{algorithmic}[1]
\Function{getrf2}{$A\in\mathbb C^{m\times n}$}
  \If{$m=1$}
    \State $\mathit{ipiv}_1\gets1$; \textbf{if} $a_{11}=0$ \textbf{then} $\mathit{info}\gets1$
  \ElsIf{$n=1$} \Comment{한 열: 피벗과 나눗셈뿐}
    \State $p\gets\arg\max_i|a_{i1}|$; $\mathit{ipiv}_1\gets p$
    \If{$a_{p1}\ne0$}
      \State $a_{11}\leftrightarrow a_{p1}$
      \State \textbf{if} $|a_{11}|\ge\mathrm{sfmin}$ \textbf{then} $a_{2:m,1}\gets a_{2:m,1}\cdot(1/a_{11})$ \textbf{else} $a_{2:m,1}\gets a_{2:m,1}/a_{11}$
    \Else\ $\mathit{info}\gets1$
    \EndIf
  \Else
    \State $n_1\gets\lfloor\min(m,n)/2\rfloor$; $n_2\gets n-n_1$
    \State \Call{getrf2}{$A_{:,1:n_1}$} \Comment{왼쪽 반 $\binom{A_{11}}{A_{21}}$}
    \State 교환 $\mathit{ipiv}_{1:n_1}$을 $A_{:,n_1+1:n}$에 적용 \Comment{\code{laswp}}
    \State $A_{12}\gets L_{11}^{-1}A_{12}$ \Comment{\code{trsm}}
    \State $A_{22}\gets A_{22}-A_{21}A_{12}$ \Comment{\code{gemm}}
    \State \Call{getrf2}{$A_{22}$}; $\mathit{ipiv}_{n_1+1:\min(m,n)}\mathrel{+}=n_1$ \Comment{$\mathit{info}$도 $n_1$만큼 옮긴다}
    \State 교환 $\mathit{ipiv}_{n_1+1:\min(m,n)}$을 $A_{:,1:n_1}$에 적용 \Comment{\code{laswp}: 왼쪽의 $L_{21}$}
  \EndIf
\EndFunction
\end{algorithmic}
\caption{재귀 LU (LAPACK \code{dgetrf2})}
\label{alg:lu-recursive}
\end{algorithm}

\cref{alg:lu-recursive}의 기저 경우는 둘이다. 한 행($m=1$)이면 할 일이 없고, 한 열($n=1$)이면 피벗을 찾아 바꾸고 나누는 것이 전부이다. 여기서 $\mathrm{sfmin}$은 역수가 넘치지 않는 가장 작은 수(LAPACK의 \code{dlamch('S')})이고 \cref{alg:lu-right}의 승수 계산과 같은 이유로 둔다. 재귀의 모든 층에서 산술의 대부분이 \code{gemm}이고, 크기가 반씩 줄어드므로 어느 층의 부분행렬은 반드시 캐시에 맞는다. 그래서 블록 폭 $b$를 정할 필요가 없고(이런 알고리즘을 캐시를 모르는(cache-oblivious) 알고리즘이라 한다) 반올림을 빼면 결과는 다시 부분 피벗의 $P$, $L$, $U$이다. 나누는 폭이 $n/2$가 아니라 $\min(m,n)/2$인 것은 가로로 넓은 행렬($n>m$)에서도 왼쪽 부분이 높이 $m$ 안에서 끝까지 소거될 수 있게 하기 위해서이다.

\cref{alg:lu-recursive}이 같은 결과를 내는 것도 같은 논증이다. 왼쪽 반을 재귀로 분해하는 동안 오른쪽 반은 아무것도 받지 않으므로, 돌아온 뒤 14행에서 그동안의 교환을 오른쪽 반에 적용하고 15--16행에서 \cref{prop:lu-blocked-same}의 증명과 같이 $U_{12}$와 Schur 보수를 만든다. 아래 오른쪽의 재귀(17행)가 새로 고른 교환은 행 $n_1+1$ 아래에서 일어나는데, 그 행들의 왼쪽 반에는 이미 승수 $L_{21}$이 적혀 있다. \cref{alg:lu-right}의 8행이 교환을 열 전체에 적용해 끝난 승수를 섞었듯이 18행이 그 교환을 왼쪽 반에도 적용해 $L_{21}$을 섞는다. 이것을 빠뜨리면 $L$의 행이 $P$와 어긋난다. \cref{fig:lu-recursion}은 재귀가 행렬을 나누는 모양이다.

\begin{figure}[t]
\centering
\begin{tikzpicture}[x=0.36cm,y=0.36cm,font=\scriptsize]
\draw (0,0) rectangle (8,10);
\draw (4,0) -- (4,10); \draw (0,6) -- (8,6);
\draw (2,0) -- (2,10); \draw (0,8) -- (4,8);
\draw (6,0) -- (6,6); \draw (4,4) -- (8,4);
\node at (1,9) {1}; \node at (1,4) {1};
\node at (3,9) {2}; \node at (3,4) {3};
\node[fill=white,inner sep=1pt] at (6,8) {4: \code{trsm}};
\node[fill=white,inner sep=1pt] at (6,3.2) {5: \code{gemm}};
\node at (5,5) {6}; \node at (7,5) {6};
\node at (5,2) {6}; \node at (7,2) {6};
\node[anchor=west,align=left,font=\small] at (9,8.5) {1, 2, 3: 왼쪽 반의 재귀\\(다시 반으로)};
\node[anchor=west,align=left,font=\small] at (9,5.5) {4: $U_{12}=L_{11}^{-1}A_{12}$};
\node[anchor=west,align=left,font=\small] at (9,3.6) {5: $A_{22}\gets A_{22}-L_{21}U_{12}$};
\node[anchor=west,align=left,font=\small] at (9,1.7) {6: 아래 오른쪽의 재귀};
\end{tikzpicture}
\caption{$10\times8$ 행렬에서 재귀 LU의 첫 두 층. 번호는 대략의 실행 순서이다. 왼쪽 반이 다시 반으로 나뉘어 두 열짜리 패널(1)에서 \code{trsm}(2)과 \code{gemm}(3)을 거쳐 내려가고, 그 뒤에 오른쪽 반 전체의 \code{trsm}(4)과 \code{gemm}(5)이 한 번씩 일어난다.}
\label{fig:lu-recursion}
\end{figure}

\subsection{What Accelerate Does}

stardust는 LU를 직접 짜지 않고 Apple Accelerate의 LAPACK을 부른다. Accelerate의 \code{getrf}는 참조 구현과 같은 인터페이스와 같은 피벗 규칙(\code{ipiv}, 같은 값이면 앞의 행)을 가지지만 내부의 블록 크기, 재귀의 모양, 행렬 곱의 커널은 공개되어 있지 않다. 측정으로 알 수 있는 것은 이렇다. \cref{fig:lu-gflops}에서 \code{dgetrf}는 $n=1536$에 초당 $2.8\times10^{11}$ 플롭으로 직접 짠 블록 LU의 세 배이고, $n\ge768$에서 여러 코어를 쓴다. 블록 없는 \code{dgetf2}는 직접 짠 루프와 같은 메모리의 벽에 있다. 그리고 다음 절에서 보듯 작은 행렬에서 \code{getrf}는 이상한 행동을 한다. 그 행동이 stardust의 구현에 규칙 두 개를 더했다(\cref{sec:lu-batches}).

\section{Using the Factors}
\label{sec:lu-using}

\subsection{Triangular Solves and Row Interchanges}

$PA=LU$를 얻으면 $Ax=b$는 $Ly=Pb$, $Ux=y$이다. 전진 대입 $y_i=(Pb)_i-\sum_{j<i}\ell_{ij}y_j$와 후진 대입 $x_i=(y_i-\sum_{j>i}u_{ij}x_j)/u_{ii}$는 각각 $n^2$ 플롭이다(나눗셈 $n$번 포함). 우변이 $k$개면 $2kn^2$이고, 분해의 $(2/3)n^3$보다 작으므로 분해를 한 번 하고 여러 우변에 다시 쓰는 것이 이득이다. LAPACK의 \code{getrs}가 이 일을 한다.

\begin{algorithm}[t]
\begin{algorithmic}[1]
\Require \code{getrf}의 결과($L$, $U$를 담은 $A$와 $\mathit{ipiv}$), $n\times k$ 우변 $B$
\Ensure $B$에 $X=A^{-1}B$를 덮어씀
\For{$i=1,\dots,n$} \Comment{\code{laswp}: 교환을 기록된 순서대로}
  \State \textbf{if} $\mathit{ipiv}_i\ne i$ \textbf{then} $B$의 행 $i$와 행 $\mathit{ipiv}_i$를 바꾼다
\EndFor
\For{$j=1,\dots,k$} \Comment{열 우선: 우변 하나씩}
  \For{$i=1,\dots,n$} \Comment{$Ly=Pb$, $\ell_{ii}=1$이라 나누지 않는다}
    \State $b_{i+1:n,j}\gets b_{i+1:n,j}-b_{ij}\,a_{i+1:n,\,i}$
  \EndFor
  \For{$i=n,\dots,1$} \Comment{$Ux=y$}
    \State $b_{ij}\gets b_{ij}/a_{ii}$; \ $b_{1:i-1,j}\gets b_{1:i-1,j}-b_{ij}\,a_{1:i-1,\,i}$
  \EndFor
\EndFor
\end{algorithmic}
\caption{인수로 풀기 (\code{getrs}의 \code{'N'} 경우; 열 지향의 대입)}
\label{alg:lu-solve}
\end{algorithm}

\cref{alg:lu-solve}의 대입은 열 지향으로 적었다. $y_i$가 정해지는 순간 그 값을 $L$의 $i$열에 곱해 아래 원소들에서 한꺼번에 빼므로 안쪽 루프가 열 우선 배열의 연속한 메모리를 지난다. 전진 대입은 열 지향이든 행 지향($y_i=b_i-\sum_j\ell_{ij}y_j$)이든 $j=1,2,\dots$의 순서로 빼므로 결과가 같고, 후진 대입은 두 꼴의 빼는 순서가 달라 끝 비트가 다를 수 있다. 행 교환은 \code{ipiv}를 앞에서부터 차례로 적용해야 한다. 앞의 교환이 뒤의 교환이 가리키는 행의 내용을 바꾸기 때문이다.

교환 목록을 순열로 바꾸는 것은 단위 순열에 같은 교환을 차례로 적용하는 것이다.

\begin{algorithm}[t]
\begin{algorithmic}[1]
\Require 교환 목록 $\mathit{ipiv}_1,\dots,\mathit{ipiv}_p$ (1부터), 행 수 $m$
\Ensure $x_{\mathit{perm}_r,:}=(LU)_{r,:}$인 순열 $\mathit{perm}$ (0부터)
\State $\mathit{perm}_r\gets r$ ($r=0,\dots,m-1$)
\For{$r=0,\dots,p-1$}
  \State $\mathit{perm}_r\leftrightarrow\mathit{perm}_{\mathit{ipiv}_{r+1}-1}$
\EndFor
\end{algorithmic}
\caption{교환 목록에서 순열로 (stardust \code{lu}의 \code{perm})}
\label{alg:lu-perm}
\end{algorithm}

\cref{alg:lu-perm}이 맞는 까닭은 $PA$의 $r$번째 행이 원래 몇 번째 행이었는지를 따라가는 것이 바로 $A$의 행 번호들에 같은 교환을 적용하는 것이기 때문이다. 교환 $r\leftrightarrow\mathit{ipiv}_r$은 행렬의 행 내용을 바꾸고 동시에 ``지금 $r$번째 자리에 있는 원래 행의 번호''도 바꾼다. 결과는 \code{x[perm]}이 \code{L @ U}와 같아지는 순열이다. scipy의 \code{lu(x, p\_indices=True)}가 돌려주는 \code{p}는 \code{x = L[p] @ U}를 만족하므로 이것의 역순열이다.

\subsection{The Inverse}

역행렬은 $AX=I$를 푸는 것이다. 우변이 $n$개이므로 \cref{alg:lu-solve}을 그대로 쓰면 풀이에 $2n^3$ 플롭, 분해까지 $(8/3)n^3$이다. LAPACK의 \code{getri}는 다른 길을 간다(\cref{alg:lu-getri}). 

\begin{theorem}
\label{thm:lu-inverse-residual}
$AX=I$를 부분 피벗 LU와 열마다의 대입으로 풀어 얻은 $\hat X$는 $|A\hat X-I|\le\gamma_{3n}P^{\mathsf T}|\hat L||\hat U||\hat X|$를 만족한다.
\end{theorem}

\begin{proof}
열 $j$에 \cref{thm:lu-solve-backward}을 쓰면 $(A+\Delta A_j)\hat x_j=e_j$, $|\Delta A_j|\le\gamma_{3n}P^{\mathsf T}|\hat L||\hat U|$이다. 따라서 $A\hat x_j-e_j=-\Delta A_j\hat x_j$이고 $|A\hat x_j-e_j|\le\gamma_{3n}P^{\mathsf T}|\hat L||\hat U||\hat x_j|$이다. 열들을 모으면 결론이다.
\end{proof}

이것은 \term{오른쪽 잔차}{right residual} $A\hat X-I$의 상한이다. \code{getri}로 얻은 역행렬은 대신 \term{왼쪽 잔차}{left residual} $\hat XA-I$가 작다. Du Croz와 Higham은 \code{getri}의 방법에 대해 $|\hat XA-I|\le c_nu|\hat X||\hat L||\hat U|$ 꼴의 상한을 보였다\cite{ducroz1992inversion}. 줄거리는 이렇다. $U^{-1}$을 열 지향으로 구하면 각 열이 삼각 풀이이므로 $|\hat X_UU-I|\lesssim c\,u|\hat X_U||U|$ 꼴의 왼쪽 상한이 나온다. 다음 단계 $XL=\hat X_U$를 행마다 풀면 $\hat X(L+\Delta L)=\hat X_U$, $|\Delta L|\lesssim c\,u|L|$이다. 둘을 곱해 $\hat XLU=\hat X_UU+\hat X\Delta LU$로 쓰고 $A=P^{\mathsf T}LU$를 넣으면 $\hat XA-I$의 각 항이 $|\hat X||L||U|$의 배수로 묶인다. 두 길의 차이는 어느 쪽에서 곱하느냐뿐이지만 쓰임새를 가른다. stardust의 정확도 시험은 \code{inv}의 결과를 $\|A\hat X-I\|_\infty/(n\epsilon(\|A\|_\infty\|\hat X\|_\infty+1))$로 재는데, \cref{thm:lu-inverse-residual}이 바로 이 양을 묶는다.

역행렬로 연립방정식을 풀어서는 안 되는 이유도 여기서 보인다. $\hat x=\hat Xb$로 풀면 $|b-A\hat x|\le|I-A\hat X||b|\lesssim\gamma_{3n}|L||U||\hat X||b|$이고, LU로 직접 풀면 $\gamma_{3n}|L||U||\hat x|$이다. $|\hat X||b|$는 $|\hat x|=|\hat Xb|$보다 훨씬 클 수 있다(상쇄가 일어나는 경우). 비용도 역행렬을 만드는 데 $2n^3$ 플롭이 들어 LU의 $(2/3)n^3$보다 세 배 비싸고, 만든 뒤 우변 하나에 드는 $2n^2$은 LU 인수로 푸는 것과 같다.

\begin{algorithm}[t]
\begin{algorithmic}[1]
\Require \code{getrf}의 결과($L\backslash U$와 $\mathit{ipiv}$)
\Ensure 그 자리에 $A^{-1}$
\State $U\gets U^{-1}$ (상삼각 역행렬, 그 자리에서) \Comment{\code{trtri}, $n^3/3$ 플롭}
\For{$j=n-1,\dots,1$} \Comment{$XL=U^{-1}$을 오른쪽 열부터}
  \State $w_{j+1:n}\gets\ell_{j+1:n,\,j}$; $\ell_{j+1:n,\,j}\gets0$
  \State $x_{:,j}\gets x_{:,j}-X_{:,\,j+1:n}\,w_{j+1:n}$ \Comment{\code{gemv}, $2n(n-j)$ 플롭}
\EndFor
\For{$j=n-1,\dots,1$} \Comment{$X\gets XP$: 열을 교환의 역순으로}
  \State \textbf{if} $\mathit{ipiv}_j\ne j$ \textbf{then} 열 $j$와 열 $\mathit{ipiv}_j$를 바꾼다
\EndFor
\end{algorithmic}
\caption{인수에서 역행렬 (LAPACK \code{getri}의 블록 없는 꼴)}
\label{alg:lu-getri}
\end{algorithm}

\cref{alg:lu-getri}의 2--5행이 맞는 까닭은 $X=A^{-1}=U^{-1}L^{-1}P$에서 $XP^{\mathsf T}L=U^{-1}$이기 때문이다. $Y=XP^{\mathsf T}$로 두고 $YL=U^{-1}$의 열 $j$를 보면 $L$이 단위 하삼각이어서 $y_j+\sum_{i>j}y_i\ell_{ij}=(U^{-1})_{:,j}$이고, 오른쪽 열 $y_{j+1},\dots,y_n$을 먼저 알면 $y_j$가 정해진다. 3행이 $L$의 열을 미리 떼어 두는 것은 그 자리에 $y_j$를 쓰기 때문이다. 플롭은 \code{trtri}의 $n^3/3$에 $\sum_j2n(n-j)\approx n^3$을 더해 $(4/3)n^3$이고 분해까지 $2n^3$이다. 단위행렬을 우변으로 둔 풀이는 $0$을 연산하므로 분해까지 $(8/3)n^3$이어서 \code{getri}가 $(2/3)n^3$ 싸다.

반면 \code{getri}는 단계가 셋이어서 작은 행렬에서는 호출 비용이 크다. \cref{tab:lu-inverse}은 두 길을 직접 잰 것이다. 이 기계에서는 $n\le1024$까지 단위행렬을 우변으로 둔 풀이가 빠르고 $n=2048$에서 \code{getri}가 앞선다. stardust가 어느 쪽을 골랐고 왜 그랬는지는 \cref{sec:lu-stardust-family}에 있다.

\begin{table}[t]
\begin{sdtable}{colspec={rrrr}, row{1}={font=\small\bfseries}}
\toprule
$n$ & \code{getrf} & \code{getri} & \code{getrs}$(I)$ \\
\midrule
$8$ & $0.50$ & $1.21$ & $0.21$ \\
$64$ & $13.3$ & $41.4$ & $11.6$ \\
$256$ & $125$ & $187$ & $179$ \\
$1024$ & $2769$ & $4954$ & $3825$ \\
$2048$ & $14445$ & $34501$ & $38641$ \\
\bottomrule
\end{sdtable}
\caption{역행렬의 두 길(Accelerate, 한 번의 호출 시간, 마이크로초, 여러 번 잰 최솟값). \code{getri}와 \code{getrs}$(I)$는 같은 \code{getrf} 결과에서 시작한다(\code{figures/48/bench.c}의 \code{inv}).}
\label{tab:lu-inverse}
\end{table}

\subsection{Determinants}

$PA=LU$에서 $\det P=(-1)^{s}$($s$는 실제로 일어난 교환의 수, 곧 $\mathit{ipiv}_i\ne i$인 $i$의 수)이고 $\det L=1$이므로
\[
\det A=(-1)^{s}\prod_{i=1}^nu_{ii}
\]
이다. 곱을 그대로 계산하면 쉽게 넘치거나 언더플로한다. $\operatorname{diag}(10^{200},10^{200},10^{-300})$의 행렬식은 $10^{100}$으로 binary64 안에 있지만 앞의 두 원소를 곱하는 순간 $10^{400}$이 넘쳐 $\infty$가 된다. 크기가 $1000$인 행렬에서 대각이 모두 $0.4$ 근처면 행렬식은 $10^{-400}$ 부근이어서 정말로 언더플로한다. 그래서 두 가지를 구분한다. 행렬식 자체가 범위를 벗어나는 경우와 중간 곱만 벗어나는 경우이다. 뒤의 경우를 정확히 처리하려면 지수를 따로 들고 다닌다.

\begin{algorithm}[t]
\begin{algorithmic}[1]
\Require $U$의 대각 $u_1,\dots,u_n$, 교환 목록 $\mathit{ipiv}$, 특이 여부
\Ensure $\det A$ (중간 곱이 범위를 벗어나도 결과가 범위 안이면 그 값)
\If{특이} \Return $0$ \EndIf
\State $d\gets1$; \textit{ok} $\gets$ \textbf{true}
\For{$i=1,\dots,n$}
  \State $d\gets d\cdot u_i$; \textbf{if} $\mathit{ipiv}_i\ne i$ \textbf{then} $d\gets-d$
  \State \textit{ok} $\gets$ \textit{ok} \textbf{and} ($d$가 유한하고 $|d|>2^{-1022}$) \Comment{복소수는 아래 본문}
\EndFor
\If{\textit{ok}} \Return $d$ \EndIf \Comment{대부분은 여기서 끝난다}
\State $d\gets1$; $e\gets0$ \Comment{$\det A=d\cdot2^{e}$로 들고 다닌다}
\For{$i=1,\dots,n$}
  \State $u_i=m_i\cdot2^{k_i}$, $|m_i|\in[1/2,1)$로 나누고 $e\gets e+k_i$ \Comment{\code{frexp}: 정확}
  \State $d\gets d\cdot m_i$; $d=m\cdot2^{k}$로 다시 나누어 $d\gets m$, $e\gets e+k$
  \State \textbf{if} $\mathit{ipiv}_i\ne i$ \textbf{then} $d\gets-d$
\EndFor
\State \Return $d\cdot2^{\max(-4000,\min(e,4000))}$ \Comment{\code{ldexp}: 여기서만 넘치거나 언더플로}
\end{algorithmic}
\caption{넘침을 피하는 행렬식 (stardust \code{det}의 \code{det\_of})}
\label{alg:lu-det}
\end{algorithm}

\cref{alg:lu-det}의 두 번째 루프는 첫 번째 루프와 정확히 같은 반올림을 한다. 2의 거듭제곱을 곱하거나 나누는 것은 결과가 정규 범위에 있는 한 정확하고, $|m_i|,|d|\in[1/2,1)$이면 곱도 $[1/4,1)$에 있어 정규 범위를 벗어나지 않기 때문이다. 그래서 첫 루프의 곱이 끝까지 가장 작은 정규 수 $2^{-1022}$보다 컸던 행렬에서는 두 루프의 가수가 비트까지 같고(곱이 바로 그 아래이면 첫 루프는 비정규 격자에서 $2^{-1022}$로 올림될 수 있어 조건을 등호 없이 둔다), 첫 루프가 중간에 넘쳤던 행렬에서는 두 번째 루프가 넘치지 않았다면 얻었을 값을 준다. 마지막 \code{ldexp}의 지수를 $\pm4000$으로 자르는 것은 정수 넘침을 막을 뿐 결과를 바꾸지 않는다. binary64의 지수 범위는 그보다 훨씬 좁아 $\pm4000$을 넘는 $e$는 어차피 $0$이나 $\infty$이기 때문이다. 복소수에서는 ``정규 범위에 있음''을 성분마다 판정한다. 각 성분이 정규 수이거나 정확한 $0$이고 둘이 함께 $0$은 아닐 때이다(실수 행렬을 복소 산술로 계산하면 허수부가 정확한 $0$이기 때문이다). 두 번째 루프에서는 실수부와 허수부 가운데 큰 쪽의 지수로 둘을 함께 나눈다. 그러면 큰 성분은 $[1/2,1)$에 오지만 작은 성분은 비정규 범위로 내려갈 수 있어서, 복소수에서 두 루프가 비트까지 같다는 앞의 논증은 작은 성분이 정규 범위에 머무는 동안에만 성립한다.

행렬식의 정확도는 두 단계로 따진다. 먼저 계산 과정이다. 계산한 인수는 $\hat L\hat U=P(A+\Delta A)$를 만족하고 대각 곱 $n-1$번의 반올림이 $1+\theta_{n-1}$을 보태므로
\[
\widehat{\det A}=\det(A+\Delta A)\,(1+\theta_{n-1}),\qquad|\theta_{n-1}|\le\gamma_{n-1},\quad|\Delta A|\le\gamma_nP^{\mathsf T}|\hat L||\hat U|
\]
이다. 다음은 행렬식이 $A$의 섭동에 얼마나 민감한가이다. $\det(A+E)=\det A\det(I+A^{-1}E)$이고 $\det(I+F)=1+\operatorname{tr}F+O(\|F\|^2)$이므로
\[
\frac{\det(A+E)-\det A}{\det A}=\operatorname{tr}(A^{-1}E)+O(\|E\|^2)=\sum_{i,j}(A^{-1})_{ji}e_{ij}+O(\|E\|^2)
\]
이다. 원소마다 $|e_{ij}|\le\epsilon|a_{ij}|$인 섭동에 대한 상대 조건수는 따라서
\[
\operatorname{cond}_{\det}(A)=\sum_{i,j}|a_{ij}(A^{-1})_{ji}|
\]
이고, stardust의 정확도 측정은 행렬식의 상대오차를 $u\cdot\operatorname{cond}_{\det}(A)$로 나누어 본다. 이 조건수는 적어도 $n$이다. $\sum_{i,j}a_{ij}(A^{-1})_{ji}=\operatorname{tr}(AA^{-1})=n$이기 때문이다.

행렬식이 정말로 범위를 벗어나는 응용(다변량 정규 분포의 로그 밀도처럼)에는 \code{slogdet}이 있다. \cref{alg:lu-slogdet}은 $\log|\det A|=\sum_i\log|u_{ii}|$와 부호를 따로 돌려준다. 부호는 실수면 $\pm1$, 복소수면 $\prod_iu_{ii}/|u_{ii}|$에 교환의 부호를 곱한 단위 복소수이고 \code{logabsdet}은 언제나 실수이다. 로그의 합은 넘치지 않는다. 각 로그의 절대오차가 $u|\log|u_{ii}||$ 정도이고 덧셈 $n-1$번이 더해지므로 합의 절대오차는 대략 $u\,n\sum_i|\log|u_{ii}||$ 이하이다. 행렬식을 $\exp$로 되돌리면 이 절대오차가 상대오차가 되므로, \code{slogdet}의 정확도는 행렬식이 범위 안일 때 \code{det}보다 나쁠 수 있다. 두 함수를 다 두는 까닭이다.

\begin{algorithm}[t]
\begin{algorithmic}[1]
\Require $U$의 대각 $u_1,\dots,u_n$, $\mathit{ipiv}$, 특이 여부
\Ensure $(\mathit{sign},\ \mathit{logabs})$
\If{특이} \Return $(0,\,-\infty)$ \EndIf
\State $\mathit{sign}\gets1$; $\mathit{logabs}\gets0$
\For{$i=1,\dots,n$}
  \State $a\gets|u_i|$ \Comment{복소수는 \code{hypot}으로 넘침 없이}
  \State $\mathit{logabs}\gets\mathit{logabs}+\log a$; $\mathit{sign}\gets\mathit{sign}\cdot(u_i/a)$
  \State \textbf{if} $\mathit{ipiv}_i\ne i$ \textbf{then} $\mathit{sign}\gets-\mathit{sign}$
\EndFor
\State \Return $(\mathit{sign},\ \mathit{logabs})$ \Comment{빈 행렬이면 $(1,0)$}
\end{algorithmic}
\caption{부호와 로그 행렬식 (stardust \code{slogdet}의 \code{sign\_logabs})}
\label{alg:lu-slogdet}
\end{algorithm}
